% --------------------------------------------------------------
%  Capítulo 3 – Especificación de Requerimientos – Producto SW
% --------------------------------------------------------------
Este capítulo detalla los límites, restricciones y los requerimientos funcionales y no funcionales del sistema \textit{DOMU}.
Los requisitos se agrupan según los módulos principales: Control de Accesos, Gestión Operativa, Finanzas, Comunicación y Soporte, incorporando además las funciones de fiscalización y participación del actor Comité.

\section{Límites del Proyecto}
En esta sección se definen los límites técnicos y operacionales del sistema, los cuales establecen el alcance previsto para su funcionamiento en condiciones reales. Estos límites permiten acotar el diseño y asegurar que la solución cumpla con lo esperado dentro de parámetros definidos y controlables.
\label{sec:limites}

\subsection{Límites de Plataforma (Hardware \& SO)}

Se detallan a continuación las condiciones mínimas requeridas en términos de dispositivos, navegadores y sistemas operativos para el correcto funcionamiento del sistema, tanto para los usuarios finales como para el entorno de ejecución del backend.

\begin{itemize}

\item La aplicación cliente se implementa como una \textbf{Progressive Web App (PWA)} accesible desde navegadores modernos compatibles con \textit{Service Workers}, tales como Google Chrome, Safari o Mozilla Firefox, en dispositivos móviles y computadores personales. La aplicación puede instalarse en la pantalla de inicio del dispositivo, ofreciendo una experiencia de uso similar a la de una aplicación nativa.

\item El soporte de \textbf{notificaciones push} depende del sistema operativo y del navegador utilizado. En dispositivos \textbf{Android}, esta funcionalidad está disponible en navegadores basados en Chromium (por ejemplo, Google Chrome) en versiones modernas del sistema operativo. En dispositivos \textbf{iOS}, el soporte para notificaciones push en aplicaciones PWA está disponible a partir de \textbf{iOS 16.4} cuando la aplicación se instala en la pantalla de inicio.

\item El lector de cédulas utilizado por el sistema debe ser compatible con códigos QR bajo el estándar \textbf{ISO/IEC 18004}~\cite{iso18004} y con códigos de barras \textbf{Code 128}. Otros formatos de codificación quedan fuera del alcance del sistema.

\item El backend del sistema se despliega sobre arquitecturas de servidor \textbf{x86\_64}. No se certifica compatibilidad con arquitecturas \textbf{ARM} en entornos productivos.

\end{itemize}

\subsection{Límites de Conectividad y Red}
Esta sección define los requisitos mínimos de red necesarios para garantizar un rendimiento adecuado y una comunicación segura entre los distintos componentes del sistema.

\begin{itemize}
    \item Requiere conexión a Internet con latencia \(<\) 150 ms y ancho de banda \(\ge\) 2 Mbps simétricos por dispositivo para cumplir RNF\_02.
    \item Todo tráfico externo usa \textbf{TLS 1.3}; versiones anteriores de TLS/SSL se rechazan.
    \item El sistema no opera en redes sin salida a Internet (modo \textit{offline} queda fuera del alcance de la versión 1.0).
\end{itemize}

\subsection{Límites de Integración Externa}
Se especifican las tecnologías y formatos de integración aceptados por el sistema, junto con las restricciones relacionadas con servicios de terceros.

\begin{itemize}
    \item Las integraciones con servicios externos se realizan mediante \textbf{API REST + JSON} sobre HTTPS. Adicionalmente, el sistema utiliza \textbf{WebSocket nativo de Javalin} para comunicación bidireccional en tiempo real (chat y notificaciones).
    \item En la version~1.0, el modulo financiero opera mediante carga de comprobantes por parte del residente y validacion manual por administracion. La integracion con pasarela de pago y la conciliacion bancaria automatica quedan fuera del alcance operativo vigente.
\end{itemize}

\subsection{Límites de Datos y Retención}
Aquí se describen las restricciones relacionadas con el almacenamiento, el volumen de datos y las políticas de retención a largo plazo.

\begin{itemize}
    \item En caso de que sea necesario un respaldo, se coordinará almacenamiento en Google Drive de archivos comprimidos respaldando la información.
    \item Base de datos soportada: \textbf{MySQL}. Otros motores (PostgreSQL, MariaDB, SQL Server, etc.) no se consideran.
    \item Tamaño lógico máximo de la base de datos: 500 GB; para volúmenes mayores se requerirá particionamiento y revisión de arquitectura.
\end{itemize}

\subsection{Límites de Escalabilidad y Uso}
Se presentan las capacidades máximas del sistema en cuanto a número de usuarios, unidades habitacionales y eventos procesados por día.

\begin{itemize}
    \item Comunidad máxima: 2.000 unidades habitacionales por instalación.
    \item Usuarios concurrentes autenticados soportados: hasta 5.000.
    \item Picos de eventos (visitas, notificaciones) limitados a 20.000 registros/día.
\end{itemize}

\section{Restricciones Técnicas}
En esta sección se presentan las restricciones técnicas que condicionan el desarrollo y despliegue del sistema. Estas restricciones están asociadas a decisiones de infraestructura, arquitectura y tecnologías obligatorias que deben cumplirse a lo largo del proyecto.

\begin{itemize}
    \item La infraestructura para el MVP será los mismos equipos en los que el código fue escrito.
    \item El sistema DOMU será desplegado en servidores proporcionados por la infraestructura tecnológica de la Universidad del Bío-Bío para efectos de pruebas y demostración del proyecto.  Por esta razón, el sistema debe adaptarse a las configuraciones, políticas de seguridad y limitaciones de acceso establecidas por dicha infraestructura, tales como restricciones en la apertura de puertos, acceso remoto y administración del servidor.
\end{itemize}

\section{Identificación de Usuario}
La identificacion de usuario en \textit{DOMU} se organiza a partir de roles funcionales que determinan el alcance operativo de cada actor dentro de la comunidad. Esta definicion es coherente con los actores vigentes del Capitulo~5 y con el modelo de control de acceso aplicado por el sistema.

\begin{table}[H]
\centering
\setlength{\tabcolsep}{4pt}
\renewcommand{\arraystretch}{1.15}
\begin{tabular}{|p{0.16\textwidth}|p{0.24\textwidth}|p{0.34\textwidth}|p{0.12\textwidth}|}
\hline
\textbf{Actor} & \textbf{Cargo(s)} & \textbf{Responsabilidad principal en el sistema} & \textbf{Privilegio} \\ \hline
Administrador & Administracion de comunidad & Configura la comunidad y gestiona usuarios, residentes, comite, unidades habitacionales, personal operativo, finanzas, incidentes, biblioteca y votaciones. & Alto \\ \hline
Conserje & Conserjeria/recepcion & Gestiona el control de visitas, el registro y retiro de encomiendas, y el reporte de incidentes operativos asociados a la atencion diaria. & Medio \\ \hline
Personal operativo & Aseo, mantencion y apoyo operativo & Ejecuta tareas asignadas por la administracion y reporta su cumplimiento con evidencia cuando corresponde. & Bajo--Medio \\ \hline
Residente & Usuario final de la comunidad & Pre-autoriza visitas, consulta publicaciones y documentos, registra incidentes, reserva espacios comunes, revisa gastos comunes, vota, usa chat y participa en el marketplace. & Bajo \\ \hline
Comite & Comite de administracion & Fiscaliza informacion financiera en modo lectura, participa en la gestion de votaciones y administra documentos asociados a actas y gobernanza comunitaria. & Medio \\ \hline
\end{tabular}
\caption{Identificacion de actores y privilegios operativos del sistema DOMU.}
\label{tab:identificacion-usuarios}
\end{table}

La autenticacion tecnica del sistema se implementa mediante credenciales personales y control de acceso basado en roles, de modo que cada usuario accede solamente a los modulos y operaciones que le corresponden segun su perfil. Esta separacion evita ambiguedades entre actores con funciones distintas, especialmente entre \textit{Conserje} y \textit{Personal operativo}, y asegura que el \textit{Comite} mantenga un rol de gobernanza y fiscalizacion consistente con los casos de uso vigentes.

\section{Requerimientos Funcionales}
A continuación, se presentan los requerimientos funcionales identificados en la etapa inicial del proyecto, los cuales se consideran fundamentales para el desarrollo de la aplicación.

{\small\setlength{\tabcolsep}{4pt}\renewcommand{\arraystretch}{1.2}
\begin{longtable}{|p{0.14\textwidth}|p{0.80\textwidth}|}
\hline
\textbf{ID} & \textbf{El sistema debe…} \\ \hline
RF\_01 & \textbf{Actor(es): CONSERJE.} Registrar la \textbf{entrada y salida de VISITAS} mediante escaneo del QR de la cedula de identidad chilena o ingreso manual validado, asociando la visita a la unidad de destino. \\ \hline
RF\_02 & \textbf{Actor(es): ADMINISTRADOR, CONSERJE, PERSONAL OPERATIVO, RESIDENTE, COMITE.} Emitir \textbf{notificaciones} asociadas a eventos relevantes de la comunidad, incluyendo visitas, encomiendas, incidentes, pagos, tareas y votaciones, segun el rol involucrado. \\ \hline
RF\_03 & \textbf{Actor(es): RESIDENTE.} Permitir al RESIDENTE \textbf{pre-autorizar visitas} mediante la aplicacion, registrando previamente la informacion necesaria del visitante para facilitar su identificacion y acceso en conserjeria. \\ \hline
RF\_04 & \textbf{Actor(es): CONSERJE, RESIDENTE, ADMINISTRADOR.} Gestionar el \textbf{registro y retiro de encomiendas}, asociando unidad de destino, evidencia opcional y cambios de estado hasta su entrega. \\ \hline
RF\_05 & \textbf{Actor(es): ADMINISTRADOR.} Gestionar \textbf{personal operativo y tareas}, permitiendo registrar personal, asignar responsables, definir prioridades y monitorear el estado de las actividades operativas. \\ \hline
RF\_06 & \textbf{Actor(es): PERSONAL OPERATIVO.} Permitir al PERSONAL OPERATIVO actualizar el estado de las tareas asignadas y declarar su cumplimiento con evidencia adjunta cuando corresponda. \\ \hline
RF\_07 & \textbf{Actor(es): ADMINISTRADOR, RESIDENTE, COMITE.} Gestionar \textbf{gastos comunes} mediante la creacion de periodos y cargos, la generacion de estados de cuenta, la carga de comprobantes por parte del RESIDENTE y la validacion manual del pago por parte del ADMINISTRADOR; el COMITE debe poder consultar esta informacion en modo lectura para fiscalizacion. \\ \hline
RF\_08 & \textbf{Actor(es): ADMINISTRADOR. [Trabajo futuro]} Incorporar la \textbf{gestion formal de morosidad y conciliacion bancaria}, incluyendo regularizacion de deuda y seguimiento administrativo integral. \\ \hline
RF\_09 & \textbf{Actor(es): RESIDENTE, ADMINISTRADOR.} Gestionar \textbf{reservas de espacios comunes} con reglas de disponibilidad, aforo, bloques horarios y restricciones aplicables por mora. \\ \hline
RF\_10 & \textbf{Actor(es): RESIDENTE.} Permitir al RESIDENTE \textbf{emitir su voto} en votaciones comunitarias abiertas, registrando la participacion de forma trazable y evitando duplicidad de sufragio. \\ \hline
RF\_11 & \textbf{Actor(es): ADMINISTRADOR, RESIDENTE.} Habilitar \textbf{canales de comunicacion comunitaria}, incluyendo un muro de publicaciones para avisos de la comunidad y acceso a interacciones comunicacionales disponibles para residentes. \\ \hline
RF\_12 & \textbf{Actor(es): RESIDENTE, CONSERJE, ADMINISTRADOR.} Permitir registrar \textbf{incidentes o reclamos operativos}, adjuntar antecedentes y mantener seguimiento de su estado por parte de administracion. \\ \hline
RF\_13 & \textbf{Actor(es): ADMINISTRADOR, COMITE.} Permitir la \textbf{gestion de votaciones comunitarias} por parte de ADMINISTRADOR y COMITE, incluyendo creacion, cierre y publicacion de resultados dentro del modulo respectivo. \\ \hline
RF\_14 & \textbf{Actor(es): ADMINISTRADOR, CONSERJE. [Trabajo futuro]} Incorporar \textbf{mantenimiento preventivo y bitacoras tecnicas especializadas}, con planificacion, alertas y evidencia de ejecucion. \\ \hline
RF\_15 & \textbf{Actor(es): ADMINISTRADOR, COMITE.} Proveer \textbf{reportes y vistas de apoyo a la gestion}, incluyendo estados financieros, resumenes de periodos y estadisticas operativas disponibles para administracion y fiscalizacion. \\ \hline
RF\_16 & \textbf{Actor(es): ADMINISTRADOR.} Permitir al ADMINISTRADOR gestionar \textbf{comunidad, roles, residentes, comite, unidades habitacionales y permisos asociados} de acuerdo con la configuracion administrativa vigente. \\ \hline
RF\_17 & \textbf{Actor(es): RESIDENTE.} Habilitar \textbf{chat en tiempo real} entre residentes mediante WebSocket, con historial de conversaciones y actualizacion inmediata de mensajes. \\ \hline
RF\_18 & \textbf{Actor(es): RESIDENTE.} Proveer un \textbf{marketplace comunitario} para publicar, actualizar y retirar articulos de compraventa o intercambio entre vecinos. \\ \hline
RF\_19 & \textbf{Actor(es): ADMINISTRADOR, COMITE, RESIDENTE.} Gestionar una \textbf{biblioteca de documentos} comunitarios, permitiendo carga de archivos por ADMINISTRADOR y COMITE segun permisos; el COMITE solo puede publicar documentos bajo la categoria \textit{actas}, mientras los usuarios autorizados pueden consultarlos. \\ \hline
RF\_20 & \textbf{Actor(es): ADMINISTRADOR, CONSERJE, PERSONAL OPERATIVO, RESIDENTE, COMITE.} Entregar \textbf{notificaciones en tiempo real} mediante WebSocket, con preferencias configurables y asociadas a eventos relevantes del sistema. \\ \hline
\caption{Requerimientos Funcionales del sistema DOMU}
\label{tab:rf}
\end{longtable}
}

\section{Requerimientos No Funcionales}
Los siguientes requerimientos no funcionales definen las condiciones de calidad que debe cumplir el sistema, más allá de sus funcionalidades. Estos aspectos aseguran que la aplicación sea confiable, eficiente, accesible y fácil de mantener.

\begin{itemize}
    \item \textbf{RNF\_01 Disponibilidad}: \(\ge\) 99{,}0\,\% mensual.
    \item \textbf{RNF\_02 Rendimiento}: 95\,\% de las peticiones REST responden \(\le\) 500\,ms bajo 50\,RPS.
    \item \textbf{RNF\_03 Accesibilidad}: contraste mínimo 4{,}5:1 y navegación por teclado.
    \item \textbf{RNF\_04 Mantenibilidad}: cobertura de tests \(\ge\) 80\,\%, API REST documentada con \textbf{Swagger UI} mediante el plugin \texttt{javalin-openapi}, accesible en \texttt{/swagger}.
\end{itemize}

\section{Interfaces Externas de Entrada}
En esta sección se describen las interfaces externas de entrada que permitirán el ingreso de datos al sistema por parte de distintos actores.

{\small
\begin{longtable}{|p{0.12\textwidth}|p{0.28\textwidth}|p{0.54\textwidth}|}
\hline
\textbf{ID} & \textbf{Nombre} & \textbf{Detalle de datos} \\ \hline
DE\_01 & Registro visita & RUN, Nombre, Unidad destino, Motivo, Foto, \textit{timestamp} \\ \hline
DE\_02 & Recepción encomienda & Código paquete, RUN residente, Foto evidencia, \textit{timestamp} \\ \hline
DE\_03 & Pago gastos comunes & RUN, Monto, Medio, Fecha, Comprobante PDF \\ \hline
DE\_04 & Reserva espacio & RUN, Espacio, \textit{timestamp reserva}, Nº asistentes \\ \hline
DE\_05 & Ticket reclamo & RUN, Tipo, Descripción, Fotos, Prioridad \\ \hline
DE\_06 & Carga de acta de asamblea & Título acta, Período, Etiqueta ``actas'', Documento PDF, \textit{timestamp} \\ \hline
\caption{Interfaces de Entrada}
\label{tab:entrada}
\end{longtable}
}

\section{Interfaces Externas de Salida}
Esta sección presenta las interfaces externas de salida que el sistema entrega a los distintos usuarios finales (reportes, notificaciones o visualizaciones).

{\small
\begin{longtable}{|p{0.12\textwidth}|p{0.28\textwidth}|p{0.46\textwidth}|p{0.12\textwidth}|}
\hline
\textbf{ID} & \textbf{Nombre} & \textbf{Datos incluidos} & \textbf{Medio} \\ \hline
IS\_01 & Reporte visitas & Fecha, RUN, Hora in/out, Unidad & PDF / Email \\ \hline
IS\_02 & Estado cuenta mensual & Periodo, Cargos, Pagos, Saldo (consulta en modo lectura para Comité) & PDF / Email \\ \hline
IS\_03 & Dashboard KPI & Indicadores clave en JSON & Web / CSV \\ \hline
IS\_04 & Notificación push & Título, Mensaje, \textit{timestamp}, \textit{deep link} & Móvil \\ \hline
IS\_05 & Acta de votación & Tema, Quórum, Resultado, \textit{hash} verificación & PDF firmado \\ \hline
\caption{Interfaces de Salida}
\label{tab:salida}
\end{longtable}
}
